\chapter[Modern Importance Sampling]{Modern Importance Sampling and Coverage-Aware Flow Training}
\label{ch:bf-modern-is}

A posterior can have several economically distinct explanations of the same
observations. A proposal fitted around one explanation may generate smooth,
well-behaved trajectories while missing the others. Importance sampling makes
this failure precise: a proposal must assign sufficient probability wherever
the posterior contributes to the expectation being estimated. A tractable
unnormalized posterior is enough to compute weights, but it does not make
unvisited regions visible. Modern methods address this difficulty by adapting
mixtures, moving through intermediate distributions, learning transport maps,
or learning stochastic dynamics. Their common purpose is to make the change
of probability measure less variable.

This chapter develops these ideas for a posterior on unconstrained parameters
$x\in\R^d$. Write its unnormalized density as $\gamma(x)$, its normalizing
constant as $Z=\int\gamma(x)\,dx\in(0,\infty)$, and its normalized density as
$p(x)=\gamma(x)/Z$. In a state-space application, $\gamma$ includes the prior,
likelihood and any parameter-chart Jacobian. If a filter supplies an approximate
likelihood, all subsequent identities concern that declared approximate
posterior. Learning a better proposal cannot remove the filter approximation.

\section{What importance weights measure}
\label{sec:bf-is-foundations}

Let $q$ be a normalized proposal with $q(x)>0$ wherever $\gamma(x)>0$, up to
sets of measure zero. Draw independent $X_i\sim q$ and define
$W_i=\gamma(X_i)/q(X_i)$. For an integrable quantity $f$, direct substitution
of the proposal density gives
\begin{align}
 \E_q[W]&=\int q(x)\frac{\gamma(x)}{q(x)}\,dx=Z,\\
 \E_q[Wf(X)]&=Z\E_p[f(X)].
 \label{eq:bf-is-change-measure}
\end{align}
Thus $\widehat Z=N^{-1}\sum_iW_i$ is unbiased. When $Z$ is unknown, the
self-normalized estimate is
\begin{equation}
 \widehat\mu=\frac{\sum_iW_if(X_i)}{\sum_iW_i},\qquad
 \mu=\E_p[f(X)].
 \label{eq:bf-is-snis}
\end{equation}
The ratio is generally biased at finite $N$. Under a law of large numbers it
is consistent. Under the appropriate second-moment condition, applying the
delta method to the numerator and denominator yields
\begin{equation}
 \sqrt N(\widehat\mu-\mu)\ \Rightarrow\
 \calN(0,\sigma_f^2),\qquad
 \sigma_f^2=\int\frac{p(x)^2}{q(x)}[f(x)-\mu]^2\,dx.
 \label{eq:bf-is-snis-variance}
\end{equation}
Indeed, the first-order ratio error is
$N^{-1}\sum_i W_i[f(X_i)-\mu]/Z$, whose mean is zero and whose variance is
$\sigma_f^2/N$. This derivation explains why a proposal useful for the posterior
mean need not be equally useful for a tail event or a rare default probability.
For a single known $f$, Cauchy--Schwarz shows that the variance-minimizing
self-normalized proposal is proportional to $p|f-\mu|$, when this expression
is meaningful. It depends on the unknown answer. Approximating $p$ itself is a
more general-purpose choice \citep[Chapter 9]{owen2013mc}.

For normalizer estimation, the variance has a direct divergence interpretation:
\begin{align}
 \frac{\Var_q(W)}{Z^2}
 &=\int\frac{p(x)^2}{q(x)}\,dx-1
 =\chi^2(p\Vert q),\\
 \frac{\Var(\widehat Z)}{Z^2}&=\frac{\chi^2(p\Vert q)}{N}.
 \label{eq:bf-is-chi-square}
\end{align}
The familiar sample diagnostic
$\widehat{\mathrm{ESS}}=(\sum_iW_i)^2/\sum_iW_i^2$ estimates a weight-concentration
quantity. With suitable moment laws,
$\widehat{\mathrm{ESS}}/N\to[1+\chi^2(p\Vert q)]^{-1}$. It is neither a count
of independent posterior draws nor a certificate that all modes were visited.
If every draw lies in one of two separated modes, nearly equal observed weights
can give an apparently excellent ESS while the omitted mode remains unknown.

\section{Why dimension and missing modes matter}

Suppose both densities factorize over their coordinates:
$p(x)=\prod_jp_j(x_j)$ and $q(x)=\prod_jq_j(x_j)$.
Tonelli's theorem then gives
\begin{equation}
 1+\chi^2(p\Vert q)
 =\prod_{j=1}^d\int\frac{p_j(x_j)^2}{q_j(x_j)}\,dx_j
 =\prod_{j=1}^d[1+\chi^2(p_j\Vert q_j)].
 \label{eq:bf-is-product}
\end{equation}
A small mismatch repeated across many coordinates can therefore produce
exponential weight variance. For $p=\calN(m,I_d)$ and $q=\calN(0,I_d)$,
$W/Z=\exp(m^\top X-\|m\|^2/2)$ and Gaussian integration gives
$\E_q[(W/Z)^2]=\exp(\|m\|^2)$. A fixed mean error per coordinate has
$\|m\|^2$ proportional to $d$.

A related sample-size analysis is due to
\citet[Theorems 1.1--1.2]{chatterjee2018sample}. They connect the threshold to $\exp\{D_{\mathrm{KL}}(p\Vert q)\}$. Their bounds involve the
concentration of $\log(p/q)$ under $p$ and the integrand, with a separate
self-normalized result. The statement is not an unconditional rule that exactly
$\exp(D_{\mathrm{KL}})$ samples suffice. It explains why the posterior's typical
set can be extremely rare under a superficially plausible proposal. The
second-moment identity above and the KL threshold answer related but distinct
questions: variance control and the onset of typical-set coverage.

Modes alone do not determine probabilities. At an isolated interior mode $m_k$
with positive-definite $H_k=-\nabla^2\log\gamma(m_k)$, the local quadratic
approximation is
\begin{equation}
 \gamma(x)\simeq\gamma(m_k)
 \exp\{-\tfrac12(x-m_k)^\top H_k(x-m_k)\},\qquad
 Z_k^{\rm Lap}=\gamma(m_k)(2\pi)^{d/2}|H_k|^{-1/2}.
 \label{eq:bf-is-laplace-mode}
\end{equation}
The determinant measures local volume. Ranking modes by height discards it.
Even the Laplace mass can fail for curved ridges, boundaries, heavy tails or
nonquadratic basins. Multistart optimization and mode merging, as developed
within the adaptive multimodal MCMC framework of
\citet[Section 2.2]{pompe2020adaptive}, provide useful candidate locations.
Finding all modes of an unrestricted nonlinear density remains a separate
problem. A finite optimizer search cannot establish exhaustiveness.

A proposal does not need one component per mode. In a product of $d/2$
double-well densities, the number of modes is $2^{d/2}$, while the formula for
the density remains short. Exploiting structure can be far more economical than
enumerating modes. This distinction is central to interpreting high-dimensional
multimodal demonstrations.

\section[Weighted kernels and adaptive mixtures]{From weighted kernels to adaptive parametric mixtures}

\citet[Sections 2.1--2.2]{west1993mixtures} uses importance-weighted kernel
mixtures to refine posterior approximations. Given weighted locations, a
normalized kernel estimate has the form
\begin{equation}
 q_{\rm next}(x)=\sum_{i=1}^N\overline W_i K_H(x-X_i),\qquad
 \overline W_i=W_i/\sum_jW_j,
 \label{eq:bf-is-kernel-mixture}
\end{equation}
where $K_H$ integrates to one. Gaussian or Student kernels produce an evaluable
proposal and can follow skewness or several basins. Every next sample must be
weighted using the density of the proposal that actually generated it.
Recycling an adaptive cloud requires corresponding multiple-importance or
sequential weighting; retroactively dividing all old draws by the latest
proposal is not automatically correct.

The classical kernel density calculation exposes a limitation. Under the
usual second-order smoothness assumptions and an isotropic bandwidth $h$,
squared integrated bias is of order $h^4$, while integrated variance is of
order $(Nh^d)^{-1}$. Balancing the two terms gives
\begin{equation}
 \mathrm{MISE}(h)\simeq C_1h^4+\frac{C_2}{Nh^d},\qquad
 h_{\rm opt}=\left(\frac{dC_2}{4C_1N}\right)^{1/(d+4)},\qquad
 \mathrm{MISE}(h_{\rm opt})=O(N^{-4/(d+4)}).
 \label{eq:bf-is-kde-rate}
\end{equation}
This is the standard fixed-dimensional density-estimation calculation, not a
theorem covering arbitrary dependent adaptive weights. Its slow exponent
explains why adding local kernels alone becomes expensive. Estimating one
integral can be easier than reconstructing a whole density; the two tasks
must not be conflated.

Consider a finite mixture $q_\eta(x)=\sum_{k=1}^K a_kq_k(x;\eta_k)$.
Forward-KL fitting minimizes $D_{\mathrm{KL}}(p\Vert q_\eta)$.
Equivalently, it maximizes $\E_p\log q_\eta$. Importance weights approximate this expectation without
requiring normalized $p$. The current responsibilities are
\begin{equation}
 r_{ik}=\frac{a_kq_k(X_i;\eta_k)}{q_\eta(X_i)}.
 \label{eq:bf-is-responsibilities}
\end{equation}
The Gaussian EM updates are
\begin{align}
 a_k^{+}&=\sum_i\overline W_i r_{ik},&
 m_k^{+}&=\frac{\sum_i\overline W_i r_{ik}X_i}{a_k^{+}},\\
 \Sigma_k^{+}&=\frac{\sum_i\overline W_i r_{ik}
 (X_i-m_k^{+})(X_i-m_k^{+})^\top}{a_k^{+}}.
 \label{eq:bf-is-weighted-em}
\end{align}
Jensen's inequality applied to the latent component label gives the EM lower
bound; maximizing it yields these weighted moments. The weights remain fixed
during each empirical EM fit. \citet[Section 2]{cappe2008mixtures} derive
adaptive updates for general mixture classes. Covariance regularization,
component removal and stopping remain numerical choices; none creates evidence
about an unvisited region.

\citet[Section 3]{raftery2010imis} develops incremental mixture importance
sampling (IMIS). It adds local Gaussian components near points with large
importance weights, estimates local covariance from neighboring points,
simulates new draws and recomputes weights against the resulting proposal
mixture. It was designed for difficult epidemiological posteriors, including
nonlinear ridges. Its residual-weight trigger identifies where the current
sampled proposal is inadequate; it cannot add a component at a region that no
search or sample has reached.

MitISEM replaces unrestricted kernel proliferation with adaptive Student-$t$
mixtures and importance-weighted EM
\citep[Sections 2 and 5; Appendix A]{hoogerheide2012mitisem}. For one
$t_\nu(m,\Sigma)$ component, the Gaussian scale-mixture representation adds a
latent precision with conditional mean
$(\nu+d)/(\nu+(x-m)^\top\Sigma^{-1}(x-m))$. This supplies robust weighted
location and scatter updates. More explicitly, with
$\tau_{ik}=(\nu_k+d)/(\nu_k+\delta_{ik})$ and
$\delta_{ik}=(X_i-m_k)^\top\Sigma_k^{-1}(X_i-m_k)$, one EM step at fixed
$\nu_k$ gives
\begin{align}
 m_k^+&=\frac{\sum_i\overline W_i r_{ik}\tau_{ik}X_i}
                 {\sum_i\overline W_i r_{ik}\tau_{ik}},\\
 \Sigma_k^+&=\frac{\sum_i\overline W_i r_{ik}\tau_{ik}
          (X_i-m_k^+)(X_i-m_k^+)^\top}{\sum_i\overline W_i r_{ik}}.
 \label{eq:bf-is-student-em}
\end{align}
Here $\Sigma_k$ is a scatter parameter, not the covariance; for $\nu_k>2$ the
covariance is $\nu_k\Sigma_k/(\nu_k-2)$. Optimizing the degrees of freedom adds
a scalar nonlinear update involving the latent precision's expected logarithm.
Large-weight points motivate extra components.
The partial MitISEM extension factors marginal and conditional distributions,
reducing the dimension of each approximation problem when that factorization
is useful. The locally inspected source is the 2011 working paper preceding
the 2012 publication; its extra sequential, permutation and partial-mixture
constructions should be distinguished from a claim to reproduce every published
experiment.

\section{Defensive and adaptive proposals}

A defensive mixture protects against a fitted proposal becoming too narrow:
\begin{equation}
 q_\lambda(x)=(1-\lambda)q_{\rm fit}(x)+\lambda q_0(x),
 \qquad 0<\lambda\le1.
 \label{eq:bf-is-defensive}
\end{equation}
If $p(x)\le Cq_0(x)$ everywhere, then $p/q_\lambda\le C/\lambda$ and
\begin{equation}
 \int\frac{p^2}{q_\lambda}\le\frac{C}{\lambda}\int p
 =\frac{C}{\lambda}.
 \label{eq:bf-is-defensive-bound}
\end{equation}
The domination assumption is the substantive protection; the word
``heavy-tailed'' by itself is insufficient. A very small $\lambda$ can still
make finite-sample exploration ineffective.

\citet[Theorems 7--9]{delyon2021safe} study a kernel estimate mixed with a safe
density and give convergence and integral-estimation results under conditions
on domination, bandwidth and the declining safe share. Their root-$N$ integral
asymptotics concern increasing sample size at fixed dimension under the stated
regularity assumptions. They do not remove the bandwidth's dimension dependence
or guarantee small finite-sample cost in high dimension.
For example, the related regularized safe scheme of \citet{korba2022mirror}
requires, among its conditions, $\lambda_n\to0$, $h_n\to0$ and
$\log n/(n h_n^d\lambda_n)\to0$. It also requires $q_0\ge c p$, boundedness
and regularity of the density and kernel, and controls such as
$(1-\eta_n)\log h_n\to0$ when powered weights are used. The factor $h_n^d$
shows explicitly where dimension remains in the convergence requirements.

\citet[Sections 2--3]{korba2022mirror} regularize importance weights using an
exponent $\eta\in(0,1]$. The population calculation is
\begin{equation}
 q(x)\left(\frac{p(x)}{q(x)}\right)^\eta
 =p(x)^\eta q(x)^{1-\eta}.
 \label{eq:bf-is-power-weights}
\end{equation}
Lower $\eta$ reduces concentration but changes the distribution fitted by those
weights. Their mirror-descent interpretation and safe kernel mixture balance
this bias against variance, with asymptotic conditions driving $\eta$ toward
one. Regularized adaptation and final estimation are different operations: an
estimate of $\E_pf$ still requires weights for $p$, unless a different estimand
is explicitly intended.

Gradient-based adaptive importance samplers (GRAMIS) use target gradients,
local curvature and interactions among proposal components to direct an
adaptive mixture \citep[Section 3 and Table 2]{elvira2023gramis}. Local
preconditioning can move a component toward a high-density region; repulsion
reduces redundant concentration. Their samples retain evaluable proposal
densities and importance correction. For equally allocated samples from $M$ proposals, their deterministic-mixture
weight denominator is $M^{-1}\sum_{j=1}^M q_j(x)$, so every sampled point is
judged against the entire current proposal mixture. The mean update combines
a backtracked preconditioned score step with repulsion:
\begin{equation}
 m_i^+=m_i+\theta_i\Sigma_i\nabla\log\gamma(m_i)
                   +\sum_{j\ne i}r_{ij}.
 \label{eq:bf-is-gramis-step}
\end{equation}
The paper specifies the scaling of $r_{ij}$ from a Poisson-field construction
and a distance-dependent interaction. Backtracking tests the score step
without the repulsion term, and the covariance uses local negative inverse
Hessian information only when positive definiteness permits it, retaining
valid prior information otherwise. Thus the method is not an unconditional
Newton step through an indefinite Hessian. The reported high-dimensional banana
examples exploit strong structure. Such examples support investigation of the
mechanism, not unrestricted mode discovery in every 50-dimensional posterior.

Gaussian-mixture variational inference (GMMVI) addresses another objective.
\citet[Sections 2--4 and Appendix A]{arenz2023gmmvi} unify natural-gradient
updates, trust regions, sample reuse and component adaptation for
$D_{\mathrm{KL}}(q\Vert p)$. A natural gradient has the form
$-F(\eta)^{-1}\nabla_\eta D_{\mathrm{KL}}$, where $F$ is the Fisher information
in the chosen variational coordinates. It accounts for how parameter changes
alter distributions. Reverse KL still weights errors under $q$, so efficient
optimization does not itself prevent missing modes. The paper's warning that
ELBO-based tuning can impair other distributional metrics is directly relevant
to evaluating a learned transport.

\section{Annealing and sequential importance sampling}
\label{sec:bf-is-ais}

Let $\gamma_0=q_0$ be normalized and let $\gamma_K=\gamma$. Intermediate
unnormalized densities may be geometric bridges
$\gamma_k=q_0^{1-\beta_k}\gamma^{\beta_k}$ for
$0=\beta_0<\cdots<\beta_K=1$. An AIS realization begins at $X_0\sim q_0$.
At stage $k$, multiply its weight by
$\gamma_k(X_{k-1})/\gamma_{k-1}(X_{k-1})$ and then apply a Markov kernel
$K_k$ invariant for $\pi_k=\gamma_k/Z_k$. With fixed kernels,
\begin{equation}
 W_{\rm AIS}=\prod_{k=1}^K
 \frac{\gamma_k(X_{k-1})}{\gamma_{k-1}(X_{k-1})},\qquad
 \E[W_{\rm AIS}f(X_K)]=\int\gamma(x)f(x)\,dx.
 \label{eq:bf-is-ais-identity}
\end{equation}
The last mutation can be omitted, in which case $f$ is evaluated at $X_{K-1}$.
To see the identity, suppose the weighted marginal before stage $k$ is
$\gamma_{k-1}(x)\,dx$. Multiplication changes it to $\gamma_k(x)\,dx$.
Integrating the next state against $K_k$ preserves this measure by invariance.
Induction from $q_0$ proves the result. An equivalent path-density derivation
uses reverse kernels; it does not require each finite chain to reach equilibrium
\citep{neal2001ais}. It does require correct invariance and support. Adapting a
kernel from the same interacting batch needs separate justification and must
not silently inherit the fixed-kernel unbiasedness statement.

Log weights accumulate as a sum, avoiding products that overflow. For the
geometric bridge the increment is
$(\beta_k-\beta_{k-1})[\log\gamma(X_{k-1})-\log q_0(X_{k-1})]$.
A small temperature change limits the immediate variation of this increment.
Good mutation is still necessary: many intermediate densities with negligible
movement can preserve a bad particle configuration at high cost.

Sequential Monte Carlo adds resampling. Given normalized weights, draw parent
indices with those probabilities, reset weights appropriately and apply
mutation. Conditional unbiasedness of the resampling operation preserves the
weighted measure in expectation. Products of the incremental normalizer
estimates give the usual SMC normalizer estimator under standard fixed-design
conditions \citep{chopin2020introduction}. Resampling trades weight degeneracy
for dependence and possible loss of ancestry. It cannot manufacture a mode that
neither the initial particles nor mutation reach.

\section{Annealed flow transport and continual training}

A learned bijection $T_k$ can move particles before they are weighted. If
$X\sim\pi_{k-1}$, change of variables gives the incremental factor
\begin{equation}
 G_k(X)=\frac{\gamma_k(T_k(X))|\det DT_k(X)|}
                    {\gamma_{k-1}(X)}.
 \label{eq:bf-is-aft-weight}
\end{equation}
Multiplying the pushed-forward law by this factor produces
$(Z_k/Z_{k-1})\pi_k$. Approximation error in the map affects weight variance;
a missing Jacobian changes the measure.

AFT learns each map by minimizing a local reverse KL. Substituting the density
of $T_k(X)$ into that KL yields, up to $\log Z_k-\log Z_{k-1}$,
\begin{equation}
 \mathcal L_k(T_k)=\E_{\pi_{k-1}}
 [\log\gamma_{k-1}(X)-\log\gamma_k(T_k(X))
                          -\log|\det DT_k(X)|].
 \label{eq:bf-is-aft-loss}
\end{equation}
\citet[Equations 7--9 and Algorithm 1]{arbel2021aft} combine transport,
importance correction, resampling and MCMC mutation. Their law of large numbers
and central limit theorem require explicit smoothness, moment and bounded-weight
conditions, with further assumptions for the CLT. Data-dependent fitted maps
cannot simply be treated as deterministic maps when asserting finite-sample
unbiasedness; the paper supplies an asymptotic analysis for its adaptive scheme.

CRAFT repeatedly runs the annealed SMC procedure and uses renewed particle
approximations to update the stage maps
\citep[Equation 5, Algorithms 1--3 and supplement]{matthews2022craft}.
Replenishing particles avoids optimizing each stage indefinitely against one
small fixed cloud. The maps learn progressively while corrected particles still
represent the annealed sequence. This is a different training organization from
fitting one whole-target flow by forward KL or FAB. Its practical advantage
depends on target cost, number of stages and mutation; neither repeated fitting
nor the transport correction eliminates finite-particle mode loss.
One complete pass therefore starts with a fresh base cloud, computes each
stage's current weighted training loss, moves points with the current map,
updates their weights using \eqref{eq:bf-is-aft-weight}, resamples when the
declared criterion calls for it, and mutates with the intermediate target.
The next pass uses the updated maps. Optimizer updates and particle transport
must follow the chosen algorithm's ordering; a map fitted using a particle
cloud cannot be silently replaced in the weight denominator by a later map.

\section{FAB: minimizing the importance-weight second moment}
\label{sec:bf-is-fab}

Let $q_\phi$ be the density of a trainable flow. FAB minimizes the alpha-two
objective
\begin{equation}
 J(\phi)=\int\frac{\gamma(x)^2}{q_\phi(x)}\,dx
 =Z^2[1+\chi^2(p\Vert q_\phi)],\qquad
 g_\phi(x)=\frac{\gamma(x)^2}{q_\phi(x)J(\phi)}.
 \label{eq:bf-fab-objective}
\end{equation}
Assume $J<\infty$, $q_\phi>0$ on the target support, and continuous
differentiability of $q_\phi$ with respect to $\phi$. Assume also an integrable
envelope for the differentiated integrand in a neighborhood of $\phi$, so that
differentiation under the integral is valid.
Since $\gamma$ does not depend on the flow parameters,
\begin{align}
 \nabla_\phi J
 &=-\int\frac{\gamma^2}{q_\phi^2}\nabla_\phi q_\phi\,dx
 =-\int\frac{\gamma^2}{q_\phi}\nabla_\phi\log q_\phi\,dx,\\
 \nabla_\phi\log J
 &=-\E_{g_\phi}[\nabla_\phi\log q_\phi(X)].
 \label{eq:bf-fab-gradient}
\end{align}
The negative density score is emphasized where $p$ is large and $q$ is too
small. This is the gradient of a chi-square objective, not ordinary forward KL,
whose gradient instead averages the same score under $p$. Minimizing $J$ or
$\log J$ has the same minimizers and descent directions, but different gradient
scales. The unknown $Z$ is irrelevant to those minimizers.

\citet[Section 3.1, Equations 4--7 and Appendix A]{midgley2023fab} approximate
expectations under $g_\phi$ by AIS. At outer iteration $t$, hold $q_t=q_{\phi_t}$
fixed during the forward sampling pass and use
\begin{align}
 \gamma_{t,\beta}(x)
 &=q_t(x)^{1-\beta}\{\gamma(x)^2/q_t(x)\}^{\beta}
 =q_t(x)^{1-2\beta}\gamma(x)^{2\beta},\\
 \log\gamma_{t,\beta}(x)
 &=(1-2\beta)\log q_t(x)+2\beta\log\gamma(x),\\
 \nabla_x\log\gamma_{t,\beta}(x)
 &=(1-2\beta)\nabla_x\log q_t(x)+2\beta\nabla_x\log\gamma(x).
 \label{eq:bf-fab-bridge}
\end{align}
At $\beta=0$ the density is $q_t$; at $\beta=1/2$ it is the unnormalized
posterior; at $\beta=1$ it is $\gamma^2/q_t$. A posterior-tempering parameter
inside a likelihood is a different quantity from this FAB bridge parameter.
They must not be identified merely because both are denoted beta in software.
The AIS log-weight increment is
\begin{equation}
 \Delta\log W_i=2\Delta\beta
 [\log\gamma(X_i)-\log q_t(X_i)].
 \label{eq:bf-fab-weight}
\end{equation}
The flow density is therefore needed at arbitrary locations reached by
mutation, not just at samples generated by the flow. HMC mutations also
require its spatial score. The author implementation offers a symmetric
Gaussian random-walk Metropolis alternative, which uses density values alone.

With AIS terminal points $X_i$ and weights $W_i$, the practical loss is
\begin{equation}
 \widehat S_t(\phi)=-\sum_i
 \operatorname{stopgrad}\!\left(\frac{W_i}{\sum_jW_j}\right)
 \log q_\phi(\operatorname{stopgrad}(X_i)).
 \label{eq:bf-fab-detached-loss}
\end{equation}
Differentiation acts only on the final log density. Differentiating through the
AIS trajectory, proposal distribution or weights computes additional terms
and is wrong relative to this estimator. With fixed valid AIS kernels,
unnormalized weights can estimate $\nabla J$ unbiasedly; the ratio in
\eqref{eq:bf-fab-detached-loss} generally introduces finite-sample bias. Reverse-KL
path-gradient estimators such as Roeder's and Vaitl's solve a different
reparameterized-gradient problem \citep{roeder2017sticking,vaitl2024fastpath};
they cannot replace this detached-density gradient unchanged.

The target $g_t$ changes when $q_t$ changes. Fitting a single frozen $g_t$ to
completion is therefore not the adaptive FAB procedure. New AIS passes and
corrected replay retain the connection to the current objective. At the ideal
solution $q=p$, $g=p$ and the auxiliary emphasis disappears.

\subsection{Replay corrections and their limits}

Suppose a stored AIS point was weighted for
$\gamma^2/q_{\rm old}$. The corresponding unnormalized target ratio is
\begin{equation}
 c_\phi(x)=\frac{\gamma(x)^2/q_\phi(x)}
                 {\gamma(x)^2/q_{\rm old}(x)}
 =\frac{q_{\rm old}(x)}{q_\phi(x)},\qquad
 \log c_\phi=\log q_{\rm old}-\log q_\phi.
 \label{eq:bf-fab-replay}
\end{equation}
Store the point, its AIS log weight and the log density used for that weight.
Sample replay indices proportional to their stored weights, then use the
minibatch loss $-B^{-1}\sum_i\operatorname{stopgrad}(c_i)\log q_\phi(X_i)$.
Before the optimizer changes $\phi$, add $\log c_i$ to each selected priority
and replace its stored log density by the currently evaluated one. Successive
corrections telescope: $(q_a/q_b)(q_b/q_c)=q_a/q_c$.

This is Algorithm 1 and Section 3.2 of \citet{midgley2023fab}. It saves expensive
posterior evaluations because old points need only a new flow density. It does
not permit uniform sampling from the buffer followed by the same correction:
that would omit the original AIS priorities. Likewise, inserting detached
corrections into a trainer that automatically renormalizes every minibatch
changes the stated replay loss.

The paper's practical sampling without replacement introduces bias. The
inspected author JAX implementation additionally caps $c_i$ in the loss, with a
configurable default of ten, while refreshing stored log priorities with the
uncapped correction. Capping bounds a stochastic update but changes its
expectation. Its activation frequency and effect require measurement. Pooled
replay, finite AIS, stale priorities and dependence are further reasons not to
call practical replay gradients exactly unbiased. A useful implementation must
preserve the distinction between the underlying objective and its finite
training procedure.

\subsection{Tail integrability is stronger than full support}

A Gaussian example demonstrates the issue. If $p=\calN(0,\sigma_p^2)$ and
$q=\calN(0,\sigma_q^2)$, then
\begin{equation}
 \frac{p(x)^2}{q(x)}=C\exp\left[-x^2
 \left(\frac{1}{\sigma_p^2}-\frac{1}{2\sigma_q^2}\right)\right].
 \label{eq:bf-fab-gaussian-tail}
\end{equation}
Its integral is finite precisely when $\sigma_q^2>\sigma_p^2/2$. Both densities
have full support even when this condition fails. A highly contracted
reverse-KL warm start may therefore create an invalid FAB auxiliary density.
Radial evaluations and weight diagnostics can expose trouble but cannot prove
a global tail bound. A defensive mixture with a proved dominating component
can guarantee a finite second moment by \eqref{eq:bf-is-defensive-bound}; it
also changes the proposal family and must be accounted for explicitly.

FAB's Appendix C studies favorable dimensional scaling under factorized
$p,q$, independent coordinate parameters and perfect intermediate transitions.
Its gradient-variance calculation motivates increasing intermediate steps with
dimension. It is not a guarantee for arbitrary dependent posteriors with local
HMC. The 32-dimensional Many-Well benchmark has $65{,}536$ modes through a
product construction; it demonstrates exploitation of structure, not an
exhaustive optimizer search over arbitrary modes.

\section{Combining FAB with IAF and NeuTra}
\label{sec:bf-fab-iaf}

FAB supplies the training procedure; an inverse autoregressive flow supplies
the density family. For stage input $z$, use
\begin{equation}
 u_j=z_j\exp\{s_j(z_{<j})\}+\mu_j(z_{<j}),\qquad
 s_j=b_j+c\tanh\{h_j(z_{<j})/c\}.
 \label{eq:bf-fab-iaf-stage}
\end{equation}
Masked neural networks ensure that $s_j$ and $\mu_j$ depend only on preceding
coordinates. The free bias $b_j$ lies outside the conditional cap. This is the
option in Hoffman's author code used by the current BayesFilter IAF; the
conditional contribution is bounded, while the complete log scale need not be.
The cap $c=2$, width sixteen and four-parameter application are local choices.
The source architecture uses ELU conditioners, author block masks, nonzero
variance-scaled kernels and coordinate reversal between three stages
\citep[Section 4.1.1 and accompanying code]{hoffman2019neutra}.

The Jacobian of each stage is triangular with diagonal $\exp(s_j)$, hence
$\log|\det Df|=\sum_js_j$. For a composition $T_\phi$ with a fixed invertible
outer affine map $a+Az$, stage log determinants add and the affine contributes
$\log|\det A|$. With standard-normal base density $\varphi_d$,
\begin{equation}
 q_\phi(x)=\varphi_d(T_\phi^{-1}(x))
 |\det DT_\phi(T_\phi^{-1}(x))|^{-1}.
 \label{eq:bf-fab-iaf-density}
\end{equation}
Sampling evaluates each autoregressive stage in parallel across its coordinates.
Evaluating the density at an arbitrary $x$ requires inversion: first recover
$z_1$, then use it to recover $z_2$, and continue. The diagonal solve is
$z_j=[u_j-\mu_j(z_{<j})]\exp[-s_j(z_{<j})]$. There is no dense determinant,
but FAB makes the inverse cost more prominent than reverse-KL training does.
Batches remain parallel; the coordinate dependency is intrinsic to this IAF
orientation. A nonlinear NAF \citep{huang2018naf} is a different scalar
transformation and is not required to combine FAB with IAF.

At fixed $x$, implicit differentiation of $T_\phi(z)=x$ gives
\begin{equation}
 \partial_\phi z=-[D_zT_\phi(z)]^{-1}\partial_\phi T_\phi(z).
 \label{eq:bf-fab-inverse-derivative}
\end{equation}
Therefore the parameter gradient of
$\log q_\phi(x)=\log\varphi_d(z)-\log|\det D_zT_\phi(z)|$ includes both
explicit parameter dependence and dependence through $z$. Holding this inverse
latent value fixed would omit terms. Autodifferentiating the native inverse
recursion computes this total derivative without materializing a dense
Jacobian. Finite differences of selected parameter directions provide an
independent check. FAB's detached $x$ is fully compatible with differentiating
$z=T_\phi^{-1}(x)$ with respect to $\phi$.

Once training is complete, freeze $\phi$ and define the NeuTra target
\begin{align}
 \widetilde p_Z(z)&=\gamma(T_\phi(z))|\det DT_\phi(z)|,\\
 \nabla_z\log\widetilde p_Z(z)
 &=DT_\phi(z)^\top\nabla_x\log\gamma(T_\phi(z))
                   +\nabla_z\log|\det DT_\phi(z)|.
 \label{eq:bf-fab-neutra-target}
\end{align}
Changing variables shows that an invariant chain for this density pushes
forward to $p$, regardless of how accurately $q_\phi$ approximates $p$.
Identity latent mass is a valid choice; mass adaptation is not a prerequisite.
The auxiliary FAB target $\gamma^2/q$ never replaces the posterior in this
formula. Better importance-weight coverage may still coexist with sharp
latent curvature, so training success does not establish efficient HMC.

For a perfectly Gaussianized posterior,
$\nabla_z\log\widetilde p_Z(z)=-z$ and
$r(z)=\log\widetilde p_Z(z)+\|z\|^2/2$ is constant. The standard post-training
check evaluates the score residual and $r$ on one thousand fresh normal-base
points, preserving finite/status failures and the residual distribution.
This checks geometry where the map generates points. Independent region
occupancy and target-weighted evidence are needed for coverage; full retained
chains and their uncertainty diagnostics are needed for posterior estimation.

\section{Stein methods and stochastic transport}

Stein variational gradient descent uses kernels in a different role from a
kernel mixture proposal. For a smooth vector field $v$, transform
$X\mapsto X+\epsilon v(X)$. Integration by parts gives
\begin{equation}
 \left.\frac{d}{d\epsilon}D_{\mathrm{KL}}(q_\epsilon\Vert p)
 \right|_{\epsilon=0}
 =-\E_q[\nabla\log p(X)^\top v(X)+\nabla\!\cdot v(X)].
 \label{eq:bf-is-stein-direction}
\end{equation}
Maximizing the negative derivative over a unit ball in a vector-valued
reproducing-kernel Hilbert space yields the particle direction
\begin{equation}
 v^*(x)\propto\frac1N\sum_i
 [k(X_i,x)\nabla\log p(X_i)+\nabla_{X_i}k(X_i,x)].
 \label{eq:bf-is-svgd}
\end{equation}
The first term attracts toward posterior mass and the second repels particles
\citep{liu2016svgd}. The evolved empirical measure does not automatically supply
an evaluable normalized density for importance weighting. Kernel bandwidth and
high-dimensional distance concentration remain material issues.

Projected SVGD exploits a likelihood-informed subspace
\citep{chen2020psvgd}. A sensitivity matrix built from likelihood scores
identifies directions where the posterior differs most strongly from the prior.
For a Gaussian prior with covariance $C$, one common form is
\begin{equation}
 H=\E_p[(\nabla\log L(X))(\nabla\log L(X))^\top],\qquad
 H\psi_j=\lambda_j C^{-1}\psi_j,
 \label{eq:bf-is-informed-subspace}
\end{equation}
where $L$ is the likelihood and the eigenvectors are orthonormal in the prior
precision metric. The associated optimal conditional-expectation approximation
averages $L$ over discarded prior coordinates. Under the functional-inequality
and prior assumptions used in that literature, its KL approximation error is
controlled by a constant times $\sum_{j>r}\lambda_j$. This statement concerns
the projected target approximation; finite particles and variational training
add errors of their own. It is not a bound on an arbitrary implementation merely
using the top $r$ eigenvectors.
The method transports coordinates in that subspace and represents the
complement using the prior conditional structure. Its error bounds depend on
assumptions and the neglected eigenvalue tail. If many directions are strongly
informed, projection has little advantage; nominal dimension alone does not
determine its suitability.

Continuous flows and controlled diffusions instead learn trajectories. For a
continuous normalizing flow $\dot X_t=v_\phi(X_t,t)$, conservation of probability
implies
\begin{equation}
 \frac{d}{dt}\log q_t(X_t)=-\nabla\!\cdot v_\phi(X_t,t).
 \label{eq:bf-is-cnf}
\end{equation}
\citet[Proposition 3.1, Equation 10 and Appendix A]{wu2025annealing} train
annealed flow segments with a reverse-KL endpoint term and an integrated
squared-velocity penalty. Substituting \eqref{eq:bf-is-cnf} gives the segment
objective, up to input-density constants,
\begin{equation}
 \E[-\log\gamma_k(X_{t_k})
      -\int_{t_{k-1}}^{t_k}\nabla\!\cdot v_\phi(X_t,t)\,dt
      +\lambda\int_{t_{k-1}}^{t_k}\|v_\phi(X_t,t)\|^2\,dt].
 \label{eq:bf-is-annealing-flow}
\end{equation}
The paper also examines discretized penalties, trace estimators, a first-order
training modification and optional Langevin refinement. These approximations
and its infinitesimal theory must be distinguished from an exact finite-step
transport. The present application retains IAF's explicit triangular
Jacobian instead of substituting this continuous-flow architecture.

Sequential controlled Langevin diffusions (SCLD) combine a learned stochastic
control with sequential importance correction and resampling
\citep[Section 2, Algorithms 1--4 and Appendix A]{chen2025scld}.
A finite path-space construction makes the correction understandable. Let
$F_k(x_k\mid x_{k-1})$ be the simulated forward transition and
$B_{k-1}(x_{k-1}\mid x_k)$ a normalized backward transition. The proposal
path density is $q_0(x_0)\prod_kF_k$ and the unnormalized target path density
is $\gamma_K(x_K)\prod_kB_{k-1}$. Their ratio factors as
\begin{equation}
 W_{0:K}=\prod_{k=1}^K
 \frac{\gamma_k(x_k)B_{k-1}(x_{k-1}\mid x_k)}
      {\gamma_{k-1}(x_{k-1})F_k(x_k\mid x_{k-1})},
 \qquad\gamma_0=q_0.
 \label{eq:bf-is-path-weight}
\end{equation}
The intermediate densities telescope; integration of the backward transitions
leaves terminal density $\gamma_K$. For nondegenerate Euler diffusion steps,
the $F_k$ are Gaussian and can be evaluated. Continuous-time versions use
Radon--Nikodym derivatives between path measures under the requisite
absolute-continuity and integrability conditions. The SCLD paper derives these
weights and trains controls using subtrajectory log-variance losses, including
off-policy sampling. An additive unknown log normalizer disappears from a
variance. For a detached reference path law $R_n$ with sufficient support, the
subtrajectory objective has the form
\begin{equation}
 \mathcal L_n(\phi)=\Var_{R_n}[\log W_{n,\phi}]
 =\E_{R_n}[(\log W_{n,\phi}-\E_{R_n}\log W_{n,\phi})^2].
 \label{eq:bf-is-scld-loss}
\end{equation}
With the reference law held fixed while differentiating, its gradient is
$2\E_{R_n}[(\log W_{n,\phi}-\E_{R_n}\log W_{n,\phi})
\nabla_\phi\log W_{n,\phi}]$. The centering term's derivative cancels because
the centered log weight has mean zero. Off-policy training separates the
sampling law from the control being updated. Section 2.3 and Theorem A.2
explain this loss and its off-policy use. Proposition 2.3 instead establishes
an exponential relative-error lower bound for a particular KL estimator on
product path measures, motivating the alternative. Zero empirical variance on
a finite buffer is not equality of complete path measures. Resampling between
segments limits concentration over a whole long path, while optional MCMC
refinement improves terminal locations.

A weighted path sampler is not automatically a tractable terminal density
$q(x)$. It can provide training samples, but distilling them into an IAF is an
additional approximation. SCLD's reported multimodal examples include a
50-dimensional, forty-component mixture. Its 1,600-dimensional log-Gaussian
Cox example is primarily evaluated through an ELBO; it is not evidence of
exhaustive unknown-mode coverage in 1,600 dimensions.

\section{Evaluation and implementation}

The broad evaluation of \citet[Sections 7--8 and Appendices D--F]{blessing2024beyond}
compares eleven variational and sampling methods across synthetic and Bayesian
targets. Its mixture examples vary dimension, including 2, 50 and 200, and show
that strong performance on an ELBO or one weight diagnostic can coexist with
poor coverage. The evidence depends on the authors' architectures, tuning and
budgets; it establishes no universal winner. It supports using several
measurements that answer different questions.

For a learned posterior proposal, correctness of density and derivatives comes
first. Next, inspect region coverage using independent information where
available, importance-weight reliability on untouched draws, and consistency
across seeds. For a NeuTra map, inspect transformed scores, numerical validity
and downstream mixing. An apparently superior extreme quantile from a short
run is descriptive until its uncertainty is assessed. A short successful test
can retain a candidate; it cannot establish a production training protocol.

The author-maintained \texttt{fab-jax} repository separates sampling, replay
and training from its bundled coupling-layer flow. Its MIT license permits
adaptation with the notice retained. BayesFilter ports those FAB components to
TensorFlow while keeping the existing IAF and batch-native posterior evaluator.
This preserves a single transport implementation and avoids a second target
implementation solely for a different autodiff backend. The port follows the
paper's normalized-weight sum in \eqref{eq:bf-fab-detached-loss}; the inspected
JAX no-replay function takes a mean of those weighted terms, adding a factor
$1/N$. The replay loss, correction cap and without-replacement choice are
recorded separately. These are explicit source differences, not interchangeable
normalizations under arbitrary optimizer settings.

\enlargethispage{\baselineskip}
Three remaining costs determine the application. First, AIS requires repeated
posterior evaluations, with scores at every intermediate HMC step or values
alone for a Metropolis mutation. Second, IAF density
evaluation at arbitrary points requires autoregressive inversion. Third, replay
initialization must provide enough distinct points to avoid fitting a tiny
static cloud. The author's recommendation is to use as few annealing and
mutation steps as can make the resulting samples usefully better than the
flow's initial samples, then calibrate on the actual problem. An example
configuration is a starting hypothesis. Source faithfulness reduces accidental
algorithm changes; target-specific measurement determines whether the resulting
training procedure is affordable and scientifically useful.
